% Source: https://artofproblemsolving.com/community/c3h1584948_how_do_you_find_the_length_of_a_line_with_an_angle
% Stable identity: aops-topic-1584948-q01
Let $M$ be the midpoint of $AB$. Triangle $OAB$ is isosceles ($OA=OB=r$), so the median $OM$ is also an altitude and an angle bisector: $OM\perp AB$ and $\angle AOM=\dfrac{\theta}{2}$.

\begin{asy}
usepackage("amsmath");
size(5.5cm);
defaultpen(fontsize(10pt));
pair O=(0,0);real r=3;pair A=r*dir(210),B=r*dir(-30),M=(A+B)/2;
draw(circle(O,r),linewidth(0.9));
draw(A--B,blue+linewidth(1.1));draw(O--A);draw(O--B);draw(O--M,dashed);
draw(arc(O,0.6,210,330),red);label("$120^\circ$",O+1.0*dir(243),red);
real s=0.22;pair u=unit(O-M),v=unit(B-M);draw(M+s*u--M+s*(u+v)--M+s*v);
dot(O);dot(A);dot(B);dot(M);
label("$O$",O,N);label("$A$",A,SW);label("$B$",B,SE);label("$M$",M,dir(250));
label("$10$",(O+A)/2,NW);label("$10$",(O+B)/2,NE);
label("\mbox{}",(min(currentpicture,user=true)+max(currentpicture,user=true))/2,black); // keeps the text after the figure black when LaTeX typesets the labels inline
\end{asy}

In the right triangle $OMA$, $AM=OA\sin\angle AOM=r\sin\dfrac{\theta}{2}$, hence
\[
AB=2AM=2r\sin\dfrac{\theta}{2}.
\]
For $r=10$ and $\theta=120^\circ$:
\[
AB=20\sin60^\circ=20\cdot\dfrac{\sqrt3}{2}=\boxed{10\sqrt3}\approx17.32 .
\]
(Compare the arc length $\dfrac{\theta}{360^\circ}\cdot2\pi r=\dfrac{20\pi}{3}\approx20.94$: the chord is the straight segment, the arc the curved path, between the same two points.)
